\section{Proofs of the resource barriers}\label{app:barrier-proofs}
\subsection[Proof of result description]{Proof of \cref{thm:description}}\label{proof:description}
\begin{proof}[Proof of \cref{thm:description}]
Let $K=2N^3$, $b_N=\lceil\log_2(4K)\rceil$, and $d_N=\lfloor N/(128b_N)\rfloor$. For sufficiently large $N$, the counting bound in Theorem~\ref{thm:construction} guarantees a flip matrix of sign-rank greater than $d_N$: indeed $\log_2(4eK)\le2b_N$, so
\[
 2Kd_N\log_2(4eK)\le N^4/16<I.
\]
Select the lexicographically first such matrix. This selection is computable. For a proposed matrix $A$, the predicate $\sr(A)\le d_N$ is the truth of the finite real sentence
\[
 \exists U\in\R^{K\times d_N},\ V\in\R^{d_N\times K}\quad
 \bigwedge_{i,j} A_{ij}\sum_{s=1}^{d_N}U_{is}V_{sj}>0.
\]
Decision of finite polynomial real sentences is the classical real-closed-field decision theorem. A constructive treatment of this decidability result is given by \citet{MC12}. Enumerate the finitely many flip tables and decide the displayed sentence for each until the first failure. The counting argument ensures termination. Small $N$ can be handled by any fixed table.

A single fixed program performs this generation and then runs the finite rational rectangle-mixture algorithm. Encode $N$ in $O(\log N)$ bits. For each fixed $\epsilon$, run at a fixed rational accuracy below $\epsilon$, whose description length is constant with respect to $N$. Conservative rational ceilings can replace the logarithmic constants in Theorem~\ref{thm:main}, making every step an ordinary finite computation. The table generation happens before the first SQ; it has no query cost. The learner therefore has $m=O(\log K)$ and constant positive tolerance, while the chosen class has $\dc>d_N=\Omega(N/\log N)$. Its program description has size $O(\log N)$. The generated table occupies far more space, which this interpretation does not charge. This contradicts the displayed polynomial implication.
\end{proof}
\subsection[Proof of result classicalbarrier]{Proof of \cref{thm:classicalbarrier}}\label{proof:classicalbarrier}
\begin{proof}[Proof of \cref{thm:classicalbarrier}]
Fix a point prior $\mu$. Apply the finite separation argument of \cref{lem:finite-minimax} to the transposed matrix, obtaining a distribution $\lambda$ over monochromatic rectangles $S_R\times T_R$ with
\[
 \forall h,\quad\E_R[\1_{T_R}(h)\mu(S_R)]\ge r,
 \qquad r=\rect(A).
\]
In the Euclidean space indexed by rectangles, put
\[
 a_h(R)=c_R\sqrt{\lambda_R}\1_{T_R}(h),\qquad
 b_x(R)=\sqrt{\lambda_R}\1_{S_R}(x).
\]
Their norms are at most one, and their inner product is $A_{hx}$ times the probability that $R$ covers $(h,x)$. These products may be zero, so this alone is not a strict representation. Add an orthogonal strict unit-vector representation scaled by $\sqrt\delta$, and scale the displayed vectors by $\sqrt{1-\delta}$. Every resulting product is strict with the required sign, and its average absolute value in each row is at least $(1-\delta)r$. Pad row norms and point norms separately in orthogonal private coordinates to make them one without changing any cross inner products. Letting $\delta\downarrow0$ proves $m_{\rm avg}(A)\ge r$.

For the VC bound, suppose $v$ points are shattered and use their uniform point prior. In any strict unit-vector representation, select one row $h_\sigma$ for each sign pattern $\sigma\in\{-1,+1\}^v$. If the minimum row-average margin is $a$, then
\[
 va\le\sum_{i=1}^v\sigma_i\langle v_{h_\sigma},u_i\rangle
 \le\left\|\sum_{i=1}^v\sigma_i u_i\right\|.
\]
Averaging the square of the last norm over independent fair signs gives $v$, because cross terms cancel. Some sign pattern has norm at most $\sqrt v$, so $a\le1/\sqrt v$. Take the supremum over representations and the infimum over priors. For $v=0$ the inequality is immediate. This proves the chain, including the strict-sign correction.
\end{proof}
\subsection[Proof of result unseen]{Proof of \cref{prop:unseen}}\label{proof:unseen}
\begin{proof}[Proof of \cref{prop:unseen}]
Condition on the sampled point indices, their observed labels, and the learner's random coins. The label at any unobserved point of $\ell_0$ remains an independent fair bit. The learner's prediction there has conditional error $1/2$. A fixed point is unobserved with probability $(1-1/N)^T$. Average over the $N$ points to get the first inequality. The second follows by induction from $(1-z)^T\ge1-Tz$ for $0\le z\le1$. Rearranging proves the sample requirement. For the final claim, $T>N/2$ and $S\ge2$ give $TS\ge N$, whereas $\dc(H_A)\le2N+1\le3N$ for every table.
\end{proof}
